%! Exponential and Logarithmic Equations and Inequalities — Unit 5, Lesson 5.6 — Group Activity (Answer Key)
\documentclass[10pt]{article}
\usepackage{precalculus-article}
\usepackage{precalculus-key}   % pulls in -boxes; do NOT also load -boxes
\newcommand{\TallMath}[1]{$\displaystyle #1\rule[-1.4em]{0pt}{3.2em}$}

\begin{document}

\pageheader{Unit 5, Lesson 5.6}{Group Activity --- Answer Key}
\namepartnerperiod

\vspace{0.08in}

\begin{scenariobox}[A Population Biologist's Equations]{plumlight}
A population biologist uses exponential and logarithmic models to study wildlife populations
and drug decay. Help her solve equations and build inverse functions that arise from her work.
\end{scenariobox}

\vspace{0.08in}

\begin{tcolorbox}[colback=white, colframe=black!40,
    title=\textbf{Tier R --- Remediate},
    fonttitle=\bfseries, arc=2mm, left=3mm, right=3mm, top=2mm, bottom=2mm]

Solve each equation. Show your steps. Then verify your solution in the original equation.

\begin{enumerate}[leftmargin=1.8em, itemsep=10pt]

\item $4^x = 64$

\vspace{6pt}
Steps: \ansline{$4^x = 4^3$ (since $64 = 4^3$), so by one-to-one property $x = 3$.}

$x = $ \ans{3} \qquad Check: $4^{3} = 64$ \checkmark? \ans{Yes}

\item $2^x = 10$

\vspace{6pt}
Steps: \ansline{$\log(2^x) = \log 10 \Rightarrow x\log 2 = 1 \Rightarrow x = 1/\log 2$}

$x = \dfrac{\log 10}{\log 2} = $ \ans{$1/\log 2$} (exact) $\approx$ \ans{3.322} (decimal)

\item $\log_3 x = 4$

\vspace{6pt}
$x = 3^4 = $ \ans{81} \qquad Check: $\log_3($ \ans{81} $) = $ \ans{4} \checkmark? \ans{Yes}

\item $\log_2(x+1) = 3$

\vspace{6pt}
$x + 1 = 2^3 = $ \ans{8} $\;\Rightarrow\; x = $ \ans{7}

Check: $\log_2($ \ans{7} $+ 1) = \log_2($ \ans{8} $) = $ \ans{3} \checkmark? \ans{Yes}

\end{enumerate}
\end{tcolorbox}

\vspace{0.08in}

\begin{tcolorbox}[colback=white, colframe=black!40,
    title=\textbf{Tier A --- Approaching Proficiency},
    fonttitle=\bfseries, arc=2mm, left=3mm, right=3mm, top=2mm, bottom=2mm]

\begin{enumerate}[leftmargin=1.8em, itemsep=10pt]

\item The biologist models drug decay: $D(t) = 200\cdot\bigl(\tfrac{1}{2}\bigr)^{t/4}$.
      Solve $D(t) = 10$. Show all steps. Give the exact answer and a decimal approximation.

\vspace{4pt}
Step 1 (divide both sides by 200): \ansline{$\bigl(\tfrac{1}{2}\bigr)^{t/4} = \tfrac{10}{200} = \tfrac{1}{20}$}

\vspace{4pt}
Step 2 (take $\log$ of both sides): \ansline{$\tfrac{t}{4}\cdot\log\tfrac{1}{2} = \log\tfrac{1}{20}$}

\vspace{4pt}
Step 3 (power rule, solve for $t$): \ansline{$t = \dfrac{4\log(1/20)}{\log(1/2)}$}

\vspace{4pt}
Exact: $t = $ \ans{$4\log(20)/\log 2$} \qquad Decimal: $t \approx $ \ans{17.29} hours

\item Solve: $\log_4 x + \log_4(x-3) = 1$. Show all steps, then check for extraneous solutions.

\vspace{4pt}
Step 1 (combine logs): \ansline{$\log_4(x(x-3)) = 1$}

\vspace{4pt}
Step 2 (convert to exponential form): \ansline{$x(x-3) = 4^1 = 4$}

\vspace{4pt}
Step 3 (solve the resulting equation): \ansline{$x^2 - 3x - 4 = 0 \Rightarrow (x-4)(x+1)=0$}

\vspace{4pt}
Possible solutions: $x = $ \ans{4} or $x = $ \ans{$-1$}

\vspace{4pt}
Check each for domain validity (argument must be $> 0$): \ansline{$x=-1$: $\log_4(-1)$ undefined; extraneous. $x=4$: $\log_4 4 = 1$ and $\log_4 1 = 0$, sum $= 1$. Valid.}

\vspace{4pt}
Final answer: $x = $ \ans{4}

\item Use the natural base identity to rewrite $3^x$ using base $e$.

\vspace{4pt}
$3^x = e^{({\color{keyred}\mathbf{\ln 3}})\cdot x} \approx e^{1.099x}$

\end{enumerate}
\end{tcolorbox}

\vspace{0.08in}

\begin{tcolorbox}[colback=white, colframe=black!40,
    title=\textbf{Tier E --- Extension},
    fonttitle=\bfseries, arc=2mm, left=3mm, right=3mm, top=2mm, bottom=2mm]

\begin{enumerate}[leftmargin=1.8em, itemsep=10pt]

\item Solve the inequality $2^x > 5^x$.

\vspace{4pt}
\ansline{Take log of both sides: $x\log 2 > x\log 5$.\\[3pt]
Subtract $x\log 5$: $x(\log 2 - \log 5) > 0$.\\[3pt]
Since $\log 2 - \log 5 = \log(2/5) < 0$, dividing reverses the inequality.\\[3pt]
$x < 0$.}

Solution: $x$ \ans{$< 0$}

\item Find the inverse of $f(x) = 4 \cdot 3^{x-2} + 1$. Show each step.

\vspace{4pt}
\ansline{Swap: $x = 4\cdot3^{y-2}+1$. Subtract 1: $x-1=4\cdot3^{y-2}$. Divide by 4: $\tfrac{x-1}{4}=3^{y-2}$.\\[3pt]
Apply $\log_3$: $\log_3\tfrac{x-1}{4} = y-2$. Add 2: $y = \log_3\tfrac{x-1}{4}+2$.}

\vspace{4pt}
$f^{-1}(x) = $ \ans{$\log_3\!\dfrac{x-1}{4} + 2$}

\vspace{6pt}
Domain of $f$: \ans{all reals} \qquad Range of $f$: \ans{$(1, \infty)$}

\vspace{4pt}
Domain of $f^{-1}$: \ans{$(1, \infty)$} \qquad Range of $f^{-1}$: \ans{all reals}

\item Find the inverse of $g(x) = \log_5(x+3) - 2$. Show each step.

\vspace{4pt}
\ansline{Swap: $x = \log_5(y+3)-2$. Add 2: $x+2=\log_5(y+3)$.\\[3pt]
Exponentiate: $5^{x+2}=y+3$. Subtract 3: $y=5^{x+2}-3$.}

\vspace{4pt}
$g^{-1}(x) = $ \ans{$5^{x+2} - 3$}

\vspace{6pt}
Domain of $g$: \ans{$(-3, \infty)$} \qquad Range of $g$: \ans{all reals}

\vspace{4pt}
Domain of $g^{-1}$: \ans{all reals} \qquad Range of $g^{-1}$: \ans{$(-3, \infty)$}

\end{enumerate}
\end{tcolorbox}

\end{document}
